\section{The integral lifting bound on every allowed cyclotomic field}

% BEGIN L82K1
\paragraph{L82K1. The level of the field and the lower cyclotomic roots.}
Fix a prime $p$, a finite unramified field $F_0/\mathbb Q_p$
and $F=F_0(\mu_{p^s})$, with $s\ge1$ for odd $p$ and $s\ge2$
for $p=2$. Use a coefficient field $L$ good at this level and
write $O$ for its integers. Use the original compatible Wach and
period records of L78W13. On the successive unramified letters
their series rings are copies of $A=O[[X]]$; put
\[
  \epsilon=1+X,\quad V=\epsilon^{p^{s-1}}-1,\qquad
  Q=\Phi_{p^s}(\epsilon),\qquad \varphi(V)=VQ .
\]
The cyclotomic Galois quotient over $F$ is $1+p^s\mathbb Z_p$:
the higher cyclotomic layers are totally ramified over $\mathbb Q_p$
and linearly disjoint from $F_0$. Use a generator $\gamma$
with cyclotomic value $u=1+p^s$. For every nonzero integer $j$,
$v_p(u^j-1)=s+v_p(j)$. Indeed a prime-to-$p$ exponent preserves
the initial valuation, and taking a $p$th power increases it by
one by the binomial expansion, also at $2$ since $s\ge2$.
Inversion handles negative $j$. In particular $u$ generates this
group topologically, by its order at every finite level.

Write $N=N(T)$ for the lattice of a crystalline representation.
At every $\zeta$ with $\zeta^{p^{s-1}}=1$ its Taylor lattice,
in $\epsilon=\zeta\exp Y$, is $D_i[[Y]]$ in invariant period
coordinates by W79H3. Since $\gamma$ fixes $\zeta$ and sends
$Y$ to $uY$, its action on the specialization of $N$ there is
the identity. Thus every entry of $\gamma(n)-n$ vanishes at
$X=\zeta-1$ for these roots. Division by the monic distinguished
polynomial $V$ leaves an integral polynomial remainder of degree
less than $p^{s-1}$. Its vanishing at the distinct roots forces
it to vanish identically. This gives $(\gamma-1)N\subset VN$.
Moreover $\gamma(V)/V\equiv u\bmod V$ by the binomial
expansion in $V$, so the action on $V^jN/V^{j+1}N$ is $u^j$.
The ratio $\gamma(Q)/Q$ is an integral unit since substitution
permutes the primitive roots, with the inverse substitution
giving divisibility in the other direction. This ratio has value
one at every lower root (there $Q=p$), hence is $1\bmod V$
by the same division. In particular for $N_Q=Q^{-h}N$ one
also has $(\gamma-1)N_Q\subset VN_Q$.
% END L82K1

% BEGIN L82K2
\paragraph{L82K2. A uniform integral bound on the Wach cokernel.}
Let $f:T_1\twoheadrightarrow T_2$ be equivariant, with both
representations crystalline and the jumps of the source in
$[-h,h]$. Strictness gives the same bounds on the quotient
and kernel. Set $H=2h$ and $w_s(b)=bs+v_p(b!)$.
On every letter W79H4 gives $Q^h\varphi N\subset N$ and
$Q^hN\subset A\varphi N$ for these lattices. We prove
\[
               V^H N_2\subset fN_1,\qquad
               p^{w_s(2H)}N_2\subset fN_1 .
\]

For a fixed reduction modulo $\lambda^d$ of the coefficient
uniformizer let $I$ be the intersection over the letters of
the annihilators of the cokernel on their reduced lattices.
As in A81R8, $X^{b_d}\in I$ for some $b_d$: after inverting
$X$ the reduced map is onto by the Witt comparison modulo
$\lambda$ and then by Nakayama for the nilpotent coefficient
ideal. This uses the tilt coordinate $X$ and the actual
surjection on coefficient lattices, independent of the level.
Frobenius on the series still substitutes $\epsilon^p$ for
$\epsilon$; the series ring modulo $\lambda^d$ is free on the
$p$ elements $\epsilon^j$, $0\le j<p$, over its Frobenius
image. This follows by grouping powers modulo $p$ after reduction
to the residue field, then lifting successively as in A81R8.
After inverting $X$ the same splitting holds and $Q$ is a unit
there since it reduces to a power of $X$.

For $b\in I$ write $Q^H b=\sum_j\epsilon^j\varphi(b_j)$.
Every $b_j$ is in $I$. Indeed for $z$ in the previous target
letter, $Q^h\varphi z$ is integral, its multiple by $b$ comes
from the source, and another factor of $Q^h$ puts it in the
span of the Frobenius of that source image. Expand its coefficients
in the displayed Frobenius basis. Uniqueness after inverting $X$
in the Frobenius-translated target basis, followed by cancellation
of $X$ on the original free modules, gives $b_j z$ in the
previous image. Intersect over the cyclic set of letters.
Thus from $V^H\sum_j\epsilon^j\varphi(a_j)\in I$ it follows
that every $V^Ha_j\in I$, using $Q^HV^H=\varphi(V^H)$.
Begin with $V^H(\epsilon-1)^{b_d}\in I$. Repeated extraction of
base-$p$ digits in the exponent of $\epsilon$ isolates each
coefficient of the polynomial $(Z-1)^{b_d}$ after any $k$ steps
with $p^k>b_d$. The highest coefficient is one, hence $V^H\in I$
at each $d$. Krull intersection on the finite $A$-module cokernel
gives the integral inclusion.

Write $N'=fN_1$ on a letter. For $H>0$ and $n\in N_2$, the
element $V^Hn$ modulo $V^HN'$ is killed both by
$\prod_{j=0}^{H-1}(\gamma-u^j)$, using the $N'$ filtration,
and by $\prod_{j=H}^{2H-1}(\gamma-u^j)$, using the $N_2$
filtration and $V^{2H}N_2\subset V^HN'$. The action on the
graded pieces of $N'$ follows by applying $f$ to L82K1.
For $0\le k<2H$, the valuation of
$\prod_{0\le j<2H,\ j\ne k}(u^k-u^j)$ is
$(2H-1)s+v_p(k!)+v_p((2H-1-k)!)\le w_s(2H-1)$.
Multiply the Lagrange polynomials of these $2H$ nodes by
$p^{w_s(2H-1)}$ to make them integral. Each is a multiple of
one of the two annihilating polynomials, and their sum is this
constant. Cancel the regular $V^H$ to get the second inclusion,
even with $w_s(2H-1)$. For $H=0$ the first inclusion itself
gives surjectivity and $w_s(0)=0$.
% END L82K2

% BEGIN L82K3
\paragraph{L82K3. Kernels and bounded pairs at this level.}
For an exact sequence $T_0\hookrightarrow T_1\twoheadrightarrow
T_2$ as in L82K2, the kernel on the Wach lattices is $N_0$.
Here is the level check in the saturation proof of A81R10.
At the height-one prime of $O[[X]]$ given by $\lambda$, the
inclusion matrix has full column rank after reduction, by the Witt
comparison and injectivity of $T_0/\lambda\to T_1/\lambda$.
Any other common prime divisor of its maximal minors would give
a point of their finite common vanishing set in the open disc
after extending coefficients, by Weierstrass preparation.
The two invertible $\gamma$ actions preserve that set. A finite
orbit gives $\epsilon^{u^j-1}=1$ for some $j>0$, so the point
has $\epsilon$ of $p$-power order. At the lower roots the map in
the standard Taylor lattices of W79H3 is the map on crystalline
spaces, which is injective even after specialization. At the
primitive $p^s$ roots W79H2--W79H3 identify the Taylor lattice
as the sum of $Y^{-a}b_aL[[Y]]$ for a filtration-adapted
basis $b_a$. Strictness lets an adapted basis of the subspace
extend to one of the middle space, giving injection at $Y=0$.
These identifications are functorial by L78W11. At roots of
higher order, the Frobenius matrices are invertible on evaluation
because $Q$ is nonzero there; passing to the preceding letter
and $\zeta^p$ preserves rank. Induction excludes every common
zero. Therefore the image equals its fraction-field span
intersected with the target, by solving in each height-one
localization and intersecting in the normal ring $O[[X]]$.
Generic dimensions identify this span as the kernel. The same
intersection inside the \'etale module holds by the faithful Witt
base change of B81L1, whose residue Laurent fields and comparison
are the same here. In particular intersection with the kernel in
$Q^{-h}N_1$, $VN_1$ and $VQ^{-h}N_1$ gives the corresponding
terms for $N_0$, by cancelling the indicated scalars.

Call a pair $(y,z)$ bounded when
$y\in Q^{-h}N$, $z\in VN$, and
$(\gamma-1)y=(\varphi-1)z$. Such a pair gives an integral
finite class. To verify the change of level in B81L2, the Witt
primitive for $y$ from LWR5 gives a triangular extension basis
whose last vector $e+b$ is fixed by the cyclotomic kernel,
has Frobenius difference $y$ and generator difference $z$.
Here $VN\subset XN$, so the generator acts trivially modulo $X$
in the enlarged basis. Its matrix and inverse preserve the
integral span, and their powers extend by continuity to the
whole procyclic quotient, including at $p=2$. Frobenius and
its inverse have only finite $Q$-poles by block inversion.
These are the hypotheses of B81L2's direct construction of an
invariant crystalline frame: L78W9 solves its coefficient
recursion at zero with denominators $1-u^j$, and L78W10
continues the frame and inverse using precisely the level-$s$
factor $Q$. In the Witt comparison $V$ evaluates to the level-one
period parameter. L78W2--L78W4 clear a power of this parameter
and evaluate the triangular basis and its inverse in crystalline periods. Hence the class represented by the
bounded pair vanishes after passage to $B_{\rm cris}$.

Conversely take an integral finite target class $[c]$. Its
extension by the trivial line is crystalline with the additional
jump zero. Apply L82K2 with $b=w_s(4h)$ to lift $p^b$ times
the trivial Wach generator to $v_2$ in the extension. Its
differences $y=(\varphi-1)v_2$ and $z=(\gamma-1)v_2$ are
bounded by the height inclusion, L82K1, and the kernel
intersection above. The Witt identity $v_2=p^be_2+b_2$ gives
the Herr class $p^b[c]$ with the sign of B81L3.
This pair lifts as an exact pair of cocycles after multiplication
by some $p^k$. Indeed B81L3's rational surjectivity proof uses
strictness and the crystalline finite dimension formula J80H1,
both valid over $F$. Lift the rational crystalline extension,
choose an integral lattice with its extra generator a sufficiently
large multiple $p^l e_1$, and let $I$ be its image in the target
extension; $I$ contains $p^l$ times that lattice of the target.
Use L82K2 on the surjection of the extended lattice onto $I$
and the Wach map for multiplication by $p^l$ into $I$. This
lifts $p^k v_2$ for some $k$ as in B81L3. Its two differences
give the required bounded pair in the source, since its projection
to the trivial line is constant.
% END L82K3

% BEGIN L82K4
\paragraph{L82K4. Integral correction of the lifting defect.}
Let $U$ be the set of jumps of the source at all embeddings
of $F$, $J=((U+\{1,\ldots,h\})\setminus\{0\})\cup
\{1,\ldots,h\}$, and set
\[
      P(B)=\prod_{j\in J}(B-u^{-j}),\qquad
      C=\prod_{i=1}^h(u^i-1).
\]
Use the kernel lattice $N=N_0$. On $N_Q/N$ the polynomial
$(\gamma-1)P(\gamma)$ is zero. This quotient is $p$-torsion
free by Weierstrass division by $Q^h$. After inverting $p$ take
Taylor expansions at its primitive $p^s$ roots. In the adapted
basis $Y^{-a}b_a$ of L82K3 the extra terms are
$Y^{-a-j}b_a$ for $1\le j\le h$, with eigenvalues $u^{-a-j}$.
The $a$ occur among the source jumps by strictness. The polynomial
therefore kills each branch and kills the integral quotient.
Also $N\cap VN_Q=VN$: multiplication by $V$ is injective on
$N_Q/N$, as is seen over $L$ at these same primitive roots,
which are disjoint from the roots of $V$.

There are two integral reductions:
\[
 CVN\subset(\gamma-1)VN+V^{h+1}N,\qquad
 (1-\varphi)V^{h+1}N=V^{h+1}N .
\]
For the first, $\prod_{i=1}^h(\gamma-u^i)$ moves $VN$ into
$V^{h+1}N$ by L82K1; subtract its value at one. For the second,
$\varphi(V^{h+1}N)\subset QV^{h+1}N$ and $Q\in(p,X)$,
a Frobenius-stable ideal. The inverse geometric series converges
in the complete series module.

Write $d(Y,Z)=(\gamma-1)Y-(\varphi-1)Z$ on
$N_Q\oplus VN$. Its values are in $VN_Q$ by L82K1 and
$\varphi(V)=VQ$. If $p^k\omega=d(Y,Z)$ for
$\omega\in VN_Q$, write $CZ=(\gamma-1)x+z_1$ with $x\in VN$
and $z_1\in V^{h+1}N$. Then
\[
 Cp^k\omega=(\gamma-1)(CY-(\varphi-1)x)-(\varphi-1)z_1 .
\]
The parenthesis is in $N_Q$ and the final term is in $N$.
Applying $P(\gamma)$ and cancelling $Cp^k$ in the torsion-free
quotient gives $P(\gamma)\omega\in N\cap VN_Q=VN$. The two
reductions show $CP(\gamma)\omega\in\operatorname{im}d$.
Subtract $C(P(\gamma)-P(1))\omega$, another $(\gamma-1)$-image
from $N_Q$, to deduce $CP(1)\omega\in\operatorname{im}d$.

Take the target pair from L82K3 and use L82K2 on the original
surjection to lift its coordinates after scaling by $p^b$:
lift $p^bQ^hy$ in $N_1$ then divide by $Q^h$, and lift
$p^bz/V$ then multiply by $V$. The differential $\omega$ of
these lifts lies in the kernel $VN_Q$. Multiplying them by
$p^k$ and subtracting $p^b$ times the exact lift in L82K3
puts $p^k\omega$ in the image of the kernel differential.
The correction just proved therefore makes these lifts into an
exact bounded pair after the additional scaling by
$p^{v_p(CP(1))}$, by subtraction of a kernel pair. L82K3
proves finiteness of the resulting class. Naturality of the Witt
primitive gives its image as $p^E[c]$, with
\[
 E=2w_s(4h)+w_s(h)+\sum_{j\in J}\bigl(s+v_p(j)\bigr).
\]
The two copies of $w_s(4h)$ come from the extension and
the original map separately. This is the original case B
$e_{\rm PB}$ for the constant lower and upper heights $h,h$.
Since $J\subset[-h,2h]\setminus\{0\}$, its sum is at most
$w_s(h)+w_s(2h)$. Legendre's formula gives
$w_s(b)\le b(s+1)$. Consequently $E\le12h(s+1)\le24hs$.
If $h=0$ the empty products equal one and $E=0$.
% END L82K4

% BEGIN L82K5
\paragraph{L82K5. The original lifting test over an arbitrary allowed subfield.}
Take every $N,h,K,L$ of the native $\mathrm{Cyc}$ test, with
$N>0$, $4\mid N$ if $p=2$, $K\subseteq F=\mathbb Q_p(\mu_q)$
and $\mu_q\subset L$, where $q=\operatorname{lcm}(p,N)$.
Write $s=\max(1,v_pN)$ and $q=p^s m$, $p\nmid m$.
The prime-to-$p$ roots generate an unramified extension: reduction
on these roots is injective and each residue root lifts uniquely
by Hensel's lemma. The field generated by $\mu_{p^s}$ is totally
ramified, by the Eisenstein polynomial $\Phi_{p^s}(1+X)$.
Thus $F=F_0(\mu_{p^s})$, where $F_0=\mathbb Q_p(\mu_m)$
is the full unramified subfield of $F$. It can be larger than
the unramified subfield of $K$. Every embedding of $F$ lands in
$L$ since $F$ is the splitting field of $X^q-1$, and the
dyadic guard gives $s\ge2$. These are the tested good-layer
and level conditions of L82K1.

Restrict the given onto lattice map to $G_F$. A crystalline
comparison basis over $K$ remains one on restriction, and the
de Rham filtered space extends its scalars to $F$. In detail in
an invariant comparison basis the new invariant coefficients are
$F_0$ for crystalline periods and $F$ for de Rham periods.
The filtered subspaces over $K$ after extension are contained
in the new filtered subspaces. For equality, $F/K$ is Galois and each new filtered
step is a Galois-stable $F$-subspace. Finite Galois descent
identifies it with the scalar extension of its invariants over
$K$, which are exactly the old step by the Galois definition of
the period filtration. Thus every label restricts to a label at $K$ with the same
jumps, and the constant bounds $[-h,h]$ apply. The surjection
and the coefficient ring $O$ are the original ones.

For an integral finite target class at $K$, restriction and
L82K4 give a finite integral lift of its $p^E$ multiple over $F$.
Corestrict that lift to $K$. Transfer preserves integrality,
commutes with the given map, and preserves finiteness because
the comparison to cohomology with $B_{\rm cris}$ coefficients
commutes with the finite coset sums. The composition with
restriction on the target is multiplication by $[F:K]$.
Removing its $p$-adic unit part costs exactly $v_p([F:K])$
in the exponent. A primitive $q$th root generates $F$ and
its minimal polynomial divides $X^q-1$, so $[F:K]\le q\le pN$.
It follows that $v_p([F:K])\le\lfloor\log_p N\rfloor+1$.
The resulting exponent is at most the native integer
\[
  e_L(p,N,h)=32(2hs+1)
       +s+\lfloor\log_p(2h)\rfloor+\lfloor\log_p N\rfloor+1 ,
\]
with the native integer base-$p$ logarithm convention
$\lfloor\log_p0\rfloor=0$. A further power of $p$ supplies
that exact integer. Equality of cohomology classes means the
requested cocycle difference is an integral coboundary.
This proves $\mathrm{LiftK}$ for the specified compatible period
records at every prime and for all tests in its original quantifier.
% END L82K5
